\section*{Capítulo \ref{ch:b3:relativistic-kinematics} --- Cinemática relativista}

\begin{solution}{exo:b3:relativistic-kinematics:1}
$\gamma = 1.005$, $1.155$, $2.294$, $7.09$, $22.4$. Estación espacial: $\beta =
\num{2.57e-5}$, $\gamma - 1 \approx \beta^2/2 = \num{3.3e-10}$; a lo
largo de seis meses ($\num{1.6e7}{}\,\unit{s}$) el reloj de la tripulación
\emph{atrasa} $\approx\qty{5}{ms}$ (solo por el efecto de velocidad; a la
baja altura de la estación, el efecto gravitatorio lo reduce, pero no lo
invierte).
\end{solution}

\begin{solution}{exo:b3:relativistic-kinematics:2}
(a) $\Delta\tau = 2d/c$. (b) En cada medio tic, el fotón recorre la
hipotenusa: $(c\Delta t/2)^2 = (v\Delta t/2)^2 + d^2$ con $d =
c\Delta\tau/2$; al despejar, $\Delta t = \Delta\tau/\sqrt{1 - v^2/c^2}$. (c)
Si $d$ cambiara con el movimiento, el paso por Pitágoras sería falso. (d)
Hágase que se crucen dos reglas idénticas, cada una con un pincel en su
extremo apuntando a la otra. Si el movimiento contrajera las longitudes
transversales, cada sistema predeciría que su propia regla pinta una marca
\emph{más allá} del extremo de la otra --- dos hechos contradictorios sobre
las mismas pinceladas en el mismo \omterm{def:b3:relativistic-kinematics:event}{suceso} de cruce. Una contradicción en un
solo \omterm{def:b3:relativistic-kinematics:event}{suceso} no está permitida: las longitudes transversales no pueden
cambiar.
\end{solution}

\begin{solution}{exo:b3:relativistic-kinematics:3}
$\gamma = 22.4$. (a) $\gamma\tau = \qty{49}{\micro s}$. (b) $v\gamma
\tau = \qty{14.8}{km}$. (c) $15/22.4 = \qty{670}{m}$. (d) Tiempo de vuelo en
el laboratorio, $\qty{50}{\micro s}$: sin relatividad,
$\eu^{-50/2.2} \approx \num{e-10}$ --- ninguno; con relatividad, el
\omterm{prop:b3:relativistic-kinematics:dilation}{tiempo propio} es $50/22.4 = \qty{2.2}{\micro s}$ y la fracción $\eu^{-1}
\approx 0.37$.
\end{solution}

\begin{solution}{exo:b3:relativistic-kinematics:4}
(a) $c\Delta t = \qty{600}{m}$, $\Delta x = \qty{300}{m}$: $\Delta s^2
= (600^2 - 300^2)\,\unit{m^2} > 0$, de género tiempo. (b) Sí: basta una
señal a $\Delta x/\Delta t = c/2$. (c) Esa misma velocidad del sistema, $v =
0.5c$; \omterm{prop:b3:relativistic-kinematics:dilation}{tiempo propio} $\Delta s/c = \sqrt{600^2 - 300^2}/c = 520/c =
\qty{1.73}{\micro s}$. (d) Ahora $c\Delta t = 300 < \Delta x = 600$: de
género espacio; ningún vínculo causal es posible; ningún sistema los lleva
al mismo lugar, pero el sistema que va a $v = c^2\Delta t/\Delta x = 0.5c$
los hace simultáneos.
\end{solution}

\begin{solution}{exo:b3:relativistic-kinematics:5}
(a) $\beta = \num{1.29e-5}$: $\gamma - 1 = \num{8.3e-11}$. (b)
$86400\,\unit{s} \times \num{8.3e-11} = \qty{7.2}{\micro s}$ al día.
(c) $c \times \qty{7.2}{\micro s} = \qty{2.2}{km}$ --- la navegación
moriría en un día. (d) El neto de $+38\,\unit{\micro s}$ significa que el
corrimiento gravitatorio al azul ($+45.7$) pesa más que el atraso cinemático
($-7.2$): a \qty{20200}{km}, estar más arriba en el potencial terrestre
acelera un reloj más de lo que la velocidad orbital lo frena.
\end{solution}

\begin{solution}{exo:b3:relativistic-kinematics:6}
(a) $1.6c/1.64 = 0.976c$. (b) $1.8c/1.81 = 0.994c$. (c) $c - u =
\dfrac{(c - u')(c - v)}{c\,(1 + u'v/c^2)}$: ambos factores positivos, luego
$u < c$. (d) La mancha es un \emph{patrón} en movimiento, no una cosa: no
viaja materia, ni energía, ni información de un punto de la cara de la Luna
al siguiente --- cada fotón fue hacia la Luna a $c$.
\end{solution}

\begin{solution}{exo:b3:relativistic-kinematics:7}
(a) Longitud de onda más corta: se acerca. $f/f_0 = 589/575 = 1.024$;
$(1+\beta)/(1-\beta) = 1.049$: $\beta = 0.024$, unos
\qty{7200}{km/s}. (b) Razón $2$: $(1+\beta)/(1-\beta) = 4$, $\beta =
0.6$. (c) Desarrollando la raíz cuadrada, $\Delta\lambda/\lambda \approx
\beta$; a \qty{600}{nm}, $\qty{1}{\angstrom} = \qty{0.1}{nm}$ da
$\beta = \num{1.7e-4}$: unos \qty{50}{km/s} por ángstrom --- el
reflejo del astrónomo. (d) El efecto Doppler transversal: dilatación
temporal pura, de orden $\beta^2$ --- medible solo con precisión
atómica.
\end{solution}

\begin{solution}{exo:b3:relativistic-kinematics:8}
(a) $20/2 = \qty{10}{m}$: cabe, justo, y las dos puertas pueden cerrarse
simultáneamente (sistema del granero). (b) En el sistema del corredor, los
\omterm{def:b3:relativistic-kinematics:event}{sucesos} ``se cierra la puerta delantera'' y ``se abre la trasera''
\emph{no} son simultáneos: la salida se abre antes de que se cierre la
entrada, y la pértiga de \qty{20}{m} enhebra un granero de \qty{5}{m} sin
quedar nunca encerrada. (c) Sistema del granero: $\Delta t = 0$ sobre $\Delta x =
\qty{10}{m}$. Sistema del corredor: $|\Delta t'| = \gamma v\Delta x/c^2 =
2 \times 0.866c \times 10/c^2 = \qty{58}{ns}$. (d) ``La pértiga está
encerrada'' afirma en silencio que dos \omterm{def:b3:relativistic-kinematics:event}{sucesos} separados (las dos puertas
cerradas) son simultáneos --- un enunciado dependiente del sistema, no un
hecho.
\end{solution}

\begin{solution}{exo:b3:relativistic-kinematics:9}
$\gamma = 1.25$ a $0.6c$. (a) Longitud en movimiento $100/1.25 = \qty{80}{m}$
a \qty{1.8e8}{m/s}: \qty{444}{ns}. (b) $240/\num{1.8e8} =
\qty{1.33}{\micro s}$. (c) $|\Delta t'| = \gamma v\Delta x/c^2 =
\qty{200}{ns}$; de $t' = \gamma(t - vx/c^2)$, la garra situada a mayor
$x$ --- la delantera --- se dispara \emph{primero} para el cohete. (d)
Estación: $\Delta s^2 = 0 - 80^2 = -6400\,\unit{m^2}$. Cohete: $\Delta
x' = \gamma\,\Delta x = \qty{100}{m}$, $c\Delta t' = \qty{60}{m}$:
$60^2 - 100^2 = -6400\,\unit{m^2}$ --- invariante.
\end{solution}

\begin{solution}{exo:b3:relativistic-kinematics:10}
(a) Terra: \qty{20}{yr}; Estela: $20/\gamma = \qty{12}{yr}$. (b) La
distancia se contrae a $8/\gamma = \qty{4.8}{ly}$, recorrida en $4.8/0.8
= \qty{6}{yr}$ de su tiempo --- coherente con los \qty{12}{yr} del viaje
de ida y vuelta. (c) \omterm{def:b3:relativistic-kinematics:event}{Simultaneidad} con el \omterm{def:b3:relativistic-kinematics:event}{suceso} del viraje $(t =
\qty{10}{yr},\ x = \qty{8}{ly})$: el sistema de ida asigna al reloj
terrestre $t - vx/c^2 = 10 - 6.4 = \qty{3.6}{yr}$; el de vuelta, $10 +
6.4 = \qty{16.4}{yr}$. Su ``ahora en la Tierra'' salta \qty{12.8}{yr}
en el quiebro; su cuenta $3.6 + 12.8 + 3.6 = \qty{20}{yr}$ coincide
exactamente con la de Terra. (d) Estela ocupa \emph{dos} \omterm{thm:b3:relativistic-kinematics:postulates}{sistemas
inerciales} unidos por una aceleración que ella siente; Terra ocupa uno. La
línea de universo recta entre los \omterm{def:b3:relativistic-kinematics:event}{sucesos} de partida y de reencuentro lleva
el mayor \omterm{prop:b3:relativistic-kinematics:dilation}{tiempo propio}; solo la de Estela está quebrada.
\end{solution}

\begin{solution}{exo:b3:relativistic-kinematics:11}
(a) Si la aplicación fuera no lineal, un movimiento uniforme no parecería
uniforme en el otro sistema, lo que distinguiría puntos de un espacio y un
tiempo homogéneos. (b) Sustituyendo una relación en la otra: $t' = \gamma t +
(1 - \gamma^2)x/\gamma v$. (c) Imponiendo $x = ct$, $x' = ct'$:
$\gamma(c - v)t = c\gamma t + c(1 - \gamma^2)ct/\gamma v$, que
se resuelve dando $\gamma^2 = 1/(1 - v^2/c^2)$. (d) La relatividad dio el
mismo $\gamma$ para los dos sentidos (ningún sistema privilegiado); el
postulado de la luz convirtió la función libre restante en el $\gamma(v)$
definido.
\end{solution}

\begin{solution}{exo:b3:relativistic-kinematics:12}
(a) Con $\cosh\varphi = \gamma$ y $\sinh\varphi = \gamma\beta$, la
transformación es exactamente la rotación hiperbólica enunciada, y
$\cosh^2 - \sinh^2 = 1$ conserva $c^2t^2 - x^2$. (b)
$\tanh(\varphi_1 + \varphi_2) = (\tanh\varphi_1 + \tanh\varphi_2)/(1 +
\tanh\varphi_1\tanh\varphi_2)$: precisamente la ley de composición ---
las velocidades no se suman, las rapideces sí. (c) $\varphi = g\tau/c$, luego
$\beta = \tanh(g\tau/c)$; $\operatorname{artanh}(0.99) = 2.65$ da
$\tau = 2.65c/g \approx \qty{8.1e7}{s} \approx 2.6$ años de tiempo de a
bordo. (d) $\varphi$ recorre todos los reales mientras $\tanh\varphi$
satura en $1$: se puede acelerar sin fin, ganando rapidez
linealmente, mientras la velocidad solo se arrastra hacia $c$.
\end{solution}

\begin{solution}{pb:b3:relativistic-kinematics:1}
\textbf{1.} Todos los muones son estrictamente idénticos y se desintegran
con un ritmo estadístico fijo: la fracción superviviente de una población
mide el \omterm{prop:b3:relativistic-kinematics:dilation}{tiempo propio} transcurrido tan bien como la aguja de un reloj.
\textbf{2.} $c\tau = \qty{3.00e8}{} \times \qty{2.2e-6}{} =
\qty{660}{m}$.
\textbf{3.} $t = \qty{15}{km}/c = \qty{50}{\micro s}$: fracción
$\eu^{-50/2.2} = \eu^{-22.7} \approx \num{1.4e-10}$.
\textbf{4.} Unas partículas que ``no pueden'' recorrer más de un kilómetro
atraviesan quince y llegan en tropel.
\textbf{5.} Una mano de $\sim\qty{100}{cm^2}$: un par de muones por
segundo; un cuerpo de sección horizontal $\sim\num{e3}\,\unit{cm^2}$:
del orden de quince entre dos latidos --- la relatividad llueve sobre
todo el mundo, siempre.
\textbf{6.} $\gamma = 1/\sqrt{1 - 0.995^2} = 10.0$.
\textbf{7.} Comparar dos recuentos de la \emph{misma} población
seleccionada a lo largo de una diferencia de altura conocida elimina toda
suposición sobre dónde y cuántos muones nacen.
\textbf{8.} $t = 1907/(0.995 \times \num{3e8}) = \qty{6.39}{\micro s}$.
\textbf{9.} $\eu^{-6.39/2.2} = 0.055$: unos $31$ por hora.
\textbf{10.} $t/\gamma = \qty{0.64}{\micro s}$.
\textbf{11.} $\eu^{-0.64/2.2} = 0.75$: unos $420$ por hora.
\textbf{12.} Medidos $408 \pm 9$: los $\approx 420$ de la relatividad
concuerdan, dentro del leve frenado de los muones en el absorbedor de los
detectores; los $31$ de la física clásica yerran en un factor trece. La
montaña decide.
\textbf{13.} $408/563 = 0.725 = \eu^{-t_{\text{propio}}/\tau}$:
$t_{\text{propio}} = \qty{0.71}{\micro s}$, luego $\gamma_{\exp} =
6.39/0.71 = 9.0$ --- de lleno en el $8.8 \pm 0.8$ de Frisch y Smith.
\textbf{14.} $1907/10 = \qty{191}{m}$.
\textbf{15.} $191/(0.995c) = \qty{0.64}{\micro s}$: la misma
desintegración del $25\%$ --- la versión del muon sobre el mismo recuento.
\textbf{16.} Un recuento de chasquidos es un conjunto de \omterm{def:b3:relativistic-kinematics:event}{sucesos} locales de
coincidencia; los \omterm{def:b3:relativistic-kinematics:event}{sucesos} y sus recuentos son invariantes --- los sistemas
pueden discrepar sobre tiempos y longitudes, nunca sobre lo que marcó un
contador.
\textbf{17.} $\Delta s^2 = (c \times \qty{6.39}{\micro s})^2 -
(\qty{1907}{m})^2 = (1917^2 - 1907^2)\,\unit{m^2} = (\qty{196}{m})^2$:
$\Delta s/c = \qty{0.65}{\micro s}$ --- el \omterm{prop:b3:relativistic-kinematics:dilation}{tiempo propio}, calculado sin
salir en ningún momento del sistema de la montaña.
\textbf{18.} Los ritmos de desintegración dependerían de la altitud por
igual para toda velocidad; en cambio, la supervivencia sigue a $\gamma$ ---
las selecciones más lentas se dilatan menos, exactamente como prescribe
$\gamma(v)$.
\textbf{19.} Vidas medias de $\gamma\tau$ con precisión de una parte en
mil (anillos de almacenamiento de muones del CERN); el anillo convierte al
muon en un viajero perpetuamente acelerado --- la ``gemela viajera'' sigue
joven aun cuando el viaje sea un viraje interminable.
\textbf{20.} $\beta = \num{8.3e-7}$: $\Delta t = \tfrac12\beta^2
\times \qty{144000}{s} = \qty{50}{ns}$.
\textbf{21.} La velocidad del avión se suma o se resta a la de la Tierra en
rotación, mientras que el efecto de la altitud (gravitatorio) es el mismo en
ambos sentidos: los vuelos hacia el este y hacia el oeste separan
limpiamente las dos contribuciones (el resultado de 1971 se ajustó a ambas).
\textbf{22.} Del \cref{exo:b3:relativistic-kinematics:5}:
\qty{7.2}{\micro s} al día.
\textbf{23.} $7 \times 7.2\,\unit{\micro s} \times c \approx
\qty{15}{km}.$
\textbf{24.} Las posiciones derivarían cientos de metros en unas horas y
kilómetros en un día --- la navegación se rompería a ojos vista la primera
mañana.
\textbf{25.} En 1963, una montaña, dos contadores y unos relojes de
\qty{2.2}{\micro s} midieron $\gamma_{\exp} = 9.0$ frente a un $10$
predicho (con el frenado, $8.8 \pm 0.8$): \omterm{prop:b3:relativistic-kinematics:dilation}{dilatación del tiempo} confirmada
con un error del $10\%$ --- y la física idéntica se corrige en silencio,
cada segundo, en todos los satélites de navegación sobre su cabeza.
\end{solution}
